\section{Thirteen simultaneous general integer certificates}
\label{sec:thirteen-general-certificates}

The power replacement adds the integer inputs
\(\alpha_{\rm pow}\) and \(\delta_{\rm pow}\).
They are the literal coordinates \texttt{power\_alpha} and
\texttt{power\_delta}. The preceding general quartic already has
eleven integer certificates, with one common reference, common
numerator and scalar point data, a global contact gauge, and the
shared affine progression of first contact arguments. Its ordinary
integer predicates accept every integer, including zero, so no extra
positivity or nonzero requirement is imposed on the two new inputs.

\begin{theorem}[The thirteen-block general transfer]
Under the retained arithmetic premises of the general integer
interface, there is one sufficiently large fixed general defining
parameter \(C_{\rm parameter}=149^N\) for which the complete power
replacement has the required undecidable instance projection on the
cofinite coding range \(x>32z\). The complete table has 979 rational
variables and 829 quadratic residuals. Their sum of squares has degree
exactly four. No effective value of the sufficient \(N\) is asserted.
\end{theorem}
\begin{proof}
First instantiate the existing general simultaneous certificate with
the thirteen inputs
\[
 b,g,t,m,w,s,h,c,\gamma,j,i,\delta_{\rm pow},\alpha_{\rm pow}.
\]
The reciprocal chart of \(i\) is unchanged. The other twelve
certificates are ordinary integer-input charts. Retain the first
eleven contact-progression indices
\[
                    0,1,2,3,4,5,6,7,8,10,11
\]
and assign indices 12 and 13 to the new inputs, respectively.
These are distinct integer indices; no argument requires consecutive
indices.

There are two contact weights per input. Replace the previous sum of
22 weights by the sum of the 26 retained and new weights. At a
consumed prime belonging to any one of the thirteen blocks, the
global-gauge soundness proof uses that block's own tests, complementary
weight sum, and pole obstruction. Its valuation and parity arguments
do not depend on the presence of precisely eleven blocks. The same
finite-family height estimates therefore give integrality of all
thirteen distinguished inputs. The numerical constants now depend
on this fixed larger family. Choose \(N\) sufficiently large for
these constants and for the original scalar completeness threshold.
This is a single choice before the coding instance is varied. The
old eleven-block sufficient threshold is not claimed to remain
unchanged.

For completeness, fix a source witness with its thirteen integer
inputs. Apply the existing ordered joint scalar preparation and the
common numerator theorem to this finite tuple. Its difference points,
normalizers, and aggregate contact shifts are now fixed. Choose one
positive integer spacing larger than the diameter of their union of
contact lists. Translating list \(\ell\) by its distinct integer
index times that spacing separates the lists, by the same interval
argument as in the original progression theorem. The spacing depends
only on these fixed lists. In particular it does not depend on the
reference index subsequently adjusted at the comparison primes.

Apply the same retained finite-family prime supplier to the resulting
fixed rational positions, enlarging finite avoidance by all thirteen
inputs' fixed data. Its all-sign order conditions give the same common
reference and common-numerator congruences. At a selected prime,
only its own weight has negative odd valuation; every other weight
is integral. The global sum and its complementary sums consequently
have the required valuations. At foreign contacts the unchanged
rational-square parity argument excludes negative odd denominators.

At the prepared finite places and at infinity, with all translations
now fixed, every weight has its old asymptotic
\(q_\nu=\kappa\lambda^{-2}+O(\lambda^{-3})\). Thus
\[
 \rho=26\kappa\lambda^{-2}+O(\lambda^{-3}),\qquad
 \frac{(\rho-q_\nu)^2}{q_\nu}
       =25^2\kappa\lambda^{-2}+O(\lambda^{-3})\longrightarrow0.
\]
The argument divides by neither 26 nor 25. The source physical tests
remain bounded on their fixed neighborhoods; detector errors tend to
zero, and the actual native-pole arguments retain their prepared
local square classes and the old positive real limit. The supplier's
finite divisibility and real-slope clauses make these conditions hold
simultaneously. Apply the same native realization independently to
each actual pole input. The affine progression reconstruction then
fills the six common coefficient equations and every remaining
private chart row.

This is the existing finite-family arithmetic proof with thirteen
members. It introduces no new prime-supplier, elliptic, height, or
native-realization premise. Its source projection is the proved Pell
power replacement. Removing the fixed finite initial range
\(x\le32z\) does not change undecidability of the selected coding
set. The coefficient census and degree assertion are verified below.
\end{proof}

The retained family proofs are
\emph{One contact gauge in a combined quadratic system},
\nolinkurl{undecidability/papers/global_gauge_quartic_checkpoint.tex};
\emph{Common point coordinates in the quartic transfer},
\nolinkurl{undecidability/papers/common_points_quartic_checkpoint.tex};
and \emph{A common arithmetic progression of contact arguments},
\nolinkurl{undecidability/general_contact_v94/affine_contact_progressions.tex}.
All changes here concern the fixed cardinality of their finite family
and its complete coefficient implementation.

\paragraph{Exact implementation and census.}
The source-only replacement removes 137 power coordinates and adds
14 coordinates, including the two new integer inputs. It removes
138 source residuals and adds nine, while replacing the definition
of \(B_5\). Applied to the preceding 982-coordinate, 854-row
general system, this gives 859 coordinates and 725 rows before the
two new integer certificates are supplied.

Each new ordinary certificate contributes 60 auxiliary coordinates
and 52 rows. Its input is already counted in the new source, so it
is not counted again. The complete census is therefore
\[
               859+2\cdot60=979,\qquad
               725+2\cdot52=829.
\]
The new source core itself has 84 coordinates and 56 residuals.

The compiler clones the old ordinary \(o\)-input chart, whose first
contact used progression index 9. Its remaining private coordinates
and rows are copied with their input and prefix renamed. Graded
rescaling of the seven common progression coefficients changes index
9 to index 12 or 13. Constant denominators are cleared and affected
rows are made primitive. This is exactly the affine substitution
\[
 \begin{aligned}
 t_\ell&=t+\ell d,\\
 u_\ell&=u+2\ell v+\ell^2w,\\
 f_\ell&=f+\ell a_1+\ell^2a_2+\ell^3a_3+\ell^4a_4.
 \end{aligned}
\]
It adds no new common progression register or equation.

The two new weights from each block are added to the three rows
where the global sum was eliminated in the final reciprocal-\(i\)
contact. These are its pole-input definition, inverse row, and gate
row. Twelve new quadratic terms implement the four new weight
summands in those three rows. No other old integer-certificate row
is altered, and no existing gauge or nonzero guard is removed.

Every residual has coefficients in the integer polynomial ring of
the fixed coefficient parameters and has degree at most two in the
counted rational variables. At least one retained residual has a
nonzero quadratic homogeneous part. The sum of the squares of those
parts cannot vanish identically over the reals, so the final degree
is exactly four.

\paragraph{Independent verification.}
The complete coefficient table and certificate are
\begin{center}\small
\nolinkurl{undecidability/pell_power_v97/power_general_m163_quadrics.json}\\
\nolinkurl{undecidability/pell_power_v97/power_general_m163_certificate.json}.
\end{center}
The independent checker imports no builder. It specifies the nine
new source residuals independently, checks every removed-coordinate
occurrence, compares every unchanged row, and audits all twelve new
global-sum terms. To check the new private certificates it starts
from the earlier \emph{unprogressed} corrected \(o\)-input chart
and directly substitutes indices 12 and 13. This is independent of
the compiler's graded scaling of its already progressed template.
It checks all 104 private rows, retained guards, parameter domains,
support, degree, counts, and predecessor and output hashes. Run
\begin{verbatim}
python undecidability/pell_power_v97/build_power_general.py
python undecidability/pell_power_v97/verify_power_general.py
\end{verbatim}
The finite checks certify the complete algebraic implementation;
the quantified arithmetic completeness and source universality are
the proofs above and the retained external coding theorem.
